\chapter{Groups 17 and 18: Taking Electrons and Avoiding Reactions}

\begin{kaichang}
Tonight, find three things at home: table salt in the kitchen, fluoride toothpaste by the sink, and povidone-iodine in the medicine cabinet. One seasons food, one protects teeth, one treats cuts. They seem unrelated.

A chemist calls them relatives. Salt contains chlorine, toothpaste fluorine, and the antiseptic iodine. All three occupy the periodic table's penultimate column, Group 17. You eat one family member daily, deliberately put another in your mouth, and apply the third to broken skin. Yet the family's original faces are frightening: chlorine gas was a First World War chemical weapon, fluorine gas ignites almost anything, and pure iodine vapor can blister skin.

Why can a new outfit turn the same group's elements from hazards into household products? Their Group 18 neighbors tell an even stranger story. Textbooks declared that this whole column would never react with anything. That belief stood for nearly seventy years, then fell to an afternoon's calculation and a flask of gas.

Guess first: did an accidental mishap overturn the iron rule that inert gases never react, or was it deliberate reasoning?
\end{kaichang}

\section{The Family's Colors and Tempers}

Group 17's formal name is the \textbf{halogens}, meaning salt-formers because their metal combinations are usually salts. Its first four members are fluorine, chlorine, bromine, and iodine. Their elemental appearances are firm handbook facts:

\begin{itemize}[itemsep=2pt]
\item \textbf{Fluorine} (\chem{F_2}): a pale yellow gas so violent that even containment is difficult.
\item \textbf{Chlorine} (\chem{Cl_2}): a yellow-green gas with a pungent disinfectant smell, related to the odor beside swimming pools.
\item \textbf{Bromine} (\chem{Br_2}): a dark reddish brown liquid, the only liquid nonmetallic element at ordinary temperature and pressure. Opening its bottle releases choking red-brown vapor.
\item \textbf{Iodine} (\chem{I_2}): a purple-black solid that becomes violet vapor on gentle heating and deposits again as sparkling flakes.
\end{itemize}

Notice the progression from gases through liquid to solid. Larger molecules pull one another more strongly, requiring higher temperatures to separate. Chapter 22 explains those intermolecular attractions; for now, note the observation.

All four eagerly seize electrons, but with different strengths. Fluorine is fiercest, reacting directly with almost every element, even gold. Its violence delayed isolation until the French chemist Moissan produced it electrolytically in 1886. Chlorine is milder but still reacts directly with iron, copper, and hydrogen. Bromine is gentler; iodine least willing, requiring heating even with iron powder.

A clear strength test makes two halogens compete for the same thing. Add chlorine water to colorless sodium bromide solution, and it soon turns orange-yellow as bromine is displaced. Chlorine takes the electrons from bromide ions:

\[\chem{Cl_2} + \chem{2NaBr} \rightarrow \chem{2NaCl} + \chem{Br_2}\]

Reverse the experiment, adding bromine water to sodium chloride, and nothing happens. The stronger takes from the weaker, never vice versa. Similarly, bromine displaces iodine, while iodine displaces neither. Electron-taking strength decreases fluorine, chlorine, bromine, iodine: a fact tested bottle by bottle, not a memorized chant.

At the far right are Group 18's helium, neon, argon, krypton, xenon, and radon, now called \textbf{noble gases}. The older name inert gases proved mistaken, as our evidence story will show. All are colorless, odorless elemental gases, famously unoccupied: avoiding bonds and reactions, traveling as single atoms. You already know them at work in luminous signs. Different gases emit different colors under current: neon's signature orange-red, argon's pale violet, and xenon's dazzling daylightlike white favored in headlights and cinema projectors. Chapter 9 explained how current kicks electrons upstairs and their descent emits characteristic colored light.

\section{One Electron Short of a Full Floor}

Why does one family take electrons while its neighbor does nothing? Zoom to outer shells.

Recall Chapter 9's electron floors and Chapter 10's table: outer-electron counts largely determine chemical personality. Halogens have 7, one short of full; noble gases have full shells, 2 for helium's first and 8 for the others.

When chlorine encounters an available electron, such as sodium's easily surrendered outer electron, accepting it into the final third-shell vacancy lowers system energy. Chapter 10 emphasized that atoms do not strive for octets or have wishes. Energy is the reason. Acceptance releases energy, and if that covers the process's other expenses, change proceeds spontaneously. Halogens usually find this profitable and seize opportunities.

After acceptance, a chlorine atom becomes a \textbf{chloride ion}, \chem{Cl^-}. Its eight-electron outer shell is low-energy and stable, no longer taking things everywhere. This unlocks the opening mystery: \textbf{chlorine gas and chlorine in salt are two electron states of one element, with completely different chemistry}. Chapter 10 mentioned it; here we dig to the bottom.

Why do noble gases stay out? Run the account. Try to \textbf{give} argon an electron: shell three is full, so the newcomer must occupy higher shell four, increasing rather than lowering energy. Nobody favors a losing transaction. Try to take an electron \textbf{from} argon: existing electrons are tightly held, and first ionization energy ranks high among elements, making theft expensive. Losing at both ends, they have avoided almost all reactions for two thousand---no, billions of---years.

Notice: unprofitable, not forbidden. No law prohibits noble-gas reactions; there is only an energy ledger. A profitable transaction could therefore draw one in. Remember this for the evidence story.

\begin{qiaoliang}
Can electron-taking strength be numbered? Yes: \textbf{electron affinity} is the energy released when a gaseous atom accepts one electron, in \ud{}{kJ/mol}. Larger values mean more release and tighter holding. Tables give chlorine about $348$, bromine $325$, iodine $295$, all in \ud{}{kJ/mol}: decreasing just like the displacement experiments.

Why weaker downward? Chapter 11's tug-of-war explains: the nucleus attracts a newcomer placed on progressively higher floors, farther away behind blankets of screening electrons. Iodine's newcomer occupies shell five, beyond much of the nucleus's reach, releasing less energy. The same account explains noble gases: full outer shells give affinities around zero or negative, so adding an electron requires payment instead of releasing energy.
\end{qiaoliang}

\section{After the Electrons Move, the World Changes}

Bring in symbols for the complete sodium--chlorine account:

\[\chem{2Na} + \chem{Cl_2} \rightarrow \chem{2NaCl}\]

Each sodium atom gives its outer electron to a chlorine atom, paired in \chem{Cl_2} molecules. The resulting \chem{Na^+} and \chem{Cl^-} ions assemble electrostatically into orderly salt crystals. Chapter 18 handles that assembly's ionic-bond and lattice-energy account. Here remember \textbf{both hazardous substances disappear after reaction}. Water-explosive sodium and lung-burning chlorine are gone, leaving white crystals for food.

This is a chemical lesson worth learning early: \textbf{an element's properties are not those of its elemental substance, nor its behavior in a particular compound}. Proton identity stays while electrons and neighbors change personality completely. Chloride is not only harmless but essential: it balances charge in stomach acid, helps sodium maintain fluid osmotic pressure, and constitutes half the salt lost through sweat.

Why is chlorine gas dangerous, conversely? Its two atoms are not yet satisfied. They share outer electrons through covalent bonds (Chapter 19), but accepting more remains energetically favorable. In moist respiratory tissues, chlorine continues taking electrons and reacting with water, producing cell-damaging substances. The same element is peaceful when satisfied and harmful when hungry: that simple.

Three legitimate family jobs show how humans turn this personality into tools.

\textbf{Disinfectants.} Waterworks add small amounts of chlorine or hypochlorites. The resulting hypochlorous acid enters bacteria and oxidizes vital molecules, killing microbes. Household 84 disinfectant contains sodium hypochlorite, \chem{NaClO}, a diluted version of the same chemistry. \textbf{Dose} is crucial. A substantial margin separates microbe-killing concentrations from clearly harmful human concentrations; waterworks keep residual chlorine in a narrow useful, safe window. A faint chlorine smell signals that window is maintained.

\textbf{Fluoride toothpaste.} The main component of tooth enamel is the mineral hydroxyapatite, with the approximate formula \chem{Ca_{10}(PO_4)_{6}(OH)_2}. Bacterial acids slowly dissolve it, beginning cavities. Fluoride replaces surface hydroxyl groups to form fluorapatite, another calcium-salt crystal but much more acid-resistant. It gives enamel tougher armor. Dose matters here too: suitable amounts prevent decay, while prolonged excess causes mottling, so children's toothpaste concentration and quantity require care.

\textbf{Halogen lamps.} In ordinary incandescent lamps, hot tungsten slowly evaporates, thinning the filament and dimming the bulb. Add a little iodine or bromine, together with roughly half protective gas such as argon, and the story changes. Evaporated tungsten reaches the cooler envelope and combines with iodine as gaseous tungsten iodide. Returning near the hot filament, it decomposes, depositing tungsten and freeing iodine for another trip. This tungsten--halogen cycle uses iodine as a free porter, extending filament life and permitting hotter, smaller, brighter headlights and stage lamps. Chapter 31 reveals a reversible reaction driven by temperature differences.

\begin{shiyan}
\textbf{Reveal Iodine with a Grain of Rice}

Materials: pharmacy povidone-iodine, the brown wound antiseptic, a little steamed bread or cooked rice, a raw potato slice, and a cotton swab.

Procedure:
\begin{enumerate}[itemsep=1pt]
\item Apply a little povidone-iodine with the swab to rice and the potato's cut surface.
\item Wait a dozen or so seconds and watch the color.
\item Apply clean water to another rice grain as a control.
\end{enumerate}

What to observe: Treated areas turn brown-yellow to blue-black or nearly black, while the water control stays unchanged. Iodine molecules enter starch's helical tunnels to make the color. Starch is a major rice and potato component, revisited in Chapter 46. Sensitive enough to detect minute iodine amounts, this reaction has served chemists and biologists for over a century.

Safety: Small skin contact with household povidone-iodine is fine, but \textbf{do not swallow it or get it in your eyes}. Discard tested food; do not eat it.
\end{shiyan}

\section{An Afternoon Overturns a Seventy-Year Myth}

\begin{zhengju}
Since argon's 1894 discovery, chemists had tried unsuccessfully to make noble gases react. By the mid-twentieth century, textbooks proclaimed their outer shells guaranteed absolute nonreactivity. Few bothered trying anymore.

In 1962, British chemist Neil Bartlett at Canada's University of British Columbia studied a vicious red solid, platinum hexafluoride \chem{PtF_6}, among the strongest electron thieves known. He found it could even take an electron from \chem{O_2}, producing an unprecedented salt.

Another investigator might simply have recorded the odd result. Bartlett took another step: he consulted tables. Removing an oxygen molecule's electron required about \ud{1175}{kJ/mol}. He recalled xenon's corresponding value, about \ud{1170}{kJ/mol}. Almost identical.

He reasoned that if \chem{PtF_6} could take oxygen's electron, and xenon's was no harder, it should take xenon's too. A full shell might not balance favorably against such a thief. This was \textbf{deliberate}, based on two public numbers and an analogy. Combining colorless xenon with dark red \chem{PtF_6} vapor yielded orange-yellow solid that same day: the first noble-gas compound, xenon hexafluoroplatinate. Later analysis refined the initial simple formula: the solid was more complicated, approximately a combination of \chem{XeF^+} and \chem{PtF_5^-}. Colleagues openly corrected the result, scientific self-correction proceeding normally.

Laboratories rushed to follow. Within months came xenon difluoride \chem{XeF_2} and tetrafluoride \chem{XeF_4}; krypton soon joined in. Seventy years of absolute certainty vanished, and the misleading name inert gases gradually yielded to noble gases.

Remember one method: Bartlett treated nonreactivity as an empirical energy-account summary, not a natural law. Such summaries can change when accounts change. He was not attacking authority; he placed oxygen's 1175 beside xenon's 1170 and asked a question anyone could ask. The difference was that he consulted the table.
\end{zhengju}

\section{How Many Satisfied Chlorides Are in a Spoonful?}

Consider electron-taking on your kitchen's scale. A level spoonful of salt weighs about 5 grams. At \chem{NaCl}'s molar mass of \ud{58.5}{g/mol}, that is $5/58.5 \approx 0.085$ mole. Each mole contains \ud{6.02\times 10^{23}}{} structural units, Chapter 5's Avogadro constant, giving

\[0.085 \times 6.02\times 10^{23} \approx 5\times 10^{22}\]

About $5\times 10^{22}$ chloride ions, each a settled chlorine atom after electron gain. Compare Earth's roughly 8 billion people, $8\times 10^{9}$. The ion count is about sixty trillion times greater. Give every person a trillion chloride ions and the spoonful still remains undistributed. Behind every ion was an energy-releasing electron transfer. Chemical change is no rare event; it has quietly occurred in astronomical numbers in your salt jar.

\section{Where This Account Fails}

Every model has boundaries. This family's history is itself a history of boundaries, so ours need special clarity.

First, \textbf{full outer shells avoiding reactions is a trend, not a law}. Xenon and krypton have full shells yet succumb to platinum hexafluoride and fluorine. Dozens of noble-gas compounds are known, mostly of xenon and krypton with strong electron thieves fluorine and oxygen. Helium, neon, and argon still form no genuine compounds under ordinary conditions. In 2016, following theoretical prediction, researchers used extreme pressure to form sodium--helium \chem{Na_2He}, but required pressures beyond Earth's center. Properly stated, noble-gas reactivity is very low and loosens downward, not absolutely zero.

Second, \textbf{decreasing fluorine--chlorine--bromine--iodine strength applies only to the same comparison}. Another ledger may show exceptions. Fluorine's electron affinity, about \ud{328}{kJ/mol}, is slightly below chlorine's \ud{348}{kJ/mol}. Fluorine is so small that an incoming electron crowds shell two, and repulsion consumes some profit. Yet fluorine remains the strongest practical electron thief because other entries compensate, as the challenge explains. Trend arrows sum several accounts; change one, and an arrow may bend.

Third, \textbf{elemental hazards do not predict an element's behavior in compounds}. We have repeatedly used this, but everyday violations justify a separate boundary. Poisonous elemental forms do not make every containing substance poisonous; essential elements do not imply edible elemental forms. Judge a substance's present electrons and neighbors, not its elemental name.

\begin{wuqu}
``Both disinfect, so mixing them doubles the effect.'' This misconception has killed people. Alkaline \chem{NaClO} disinfectant and acidic toilet cleaners can react when mixed, releasing chlorine: acidified hypochlorite reacts with chloride and regenerates \chem{Cl_2}. Small, poorly ventilated bathrooms magnify the danger, and news repeatedly reports hospitalizations after cleaner mixing. No recipes or proportions are needed, only this principle: \textbf{chemical compatibility follows particle-level accounts, not shelf categories}. Two cleaning labels do not promise peaceful molecules. Use one household cleaner at a time and rinse thoroughly before changing products.
\end{wuqu}

\section{Back to the Three Household Products}

We can now settle the opening examples.

Salt's chlorine is chloride, electron-satisfied with a full shell, low energy, and quiet personality, stacked neatly with sodium ions. It no longer seizes anything, so you eat it daily and your body needs it.

Toothpaste's fluorine is likewise fluoride after electron gain, working at the tooth surface to make more acid-resistant fluorapatite. It no longer corrodes things: its electron-taking power was spent during manufacture.

Povidone-iodine contains elemental iodine in solution, not yet electron-satisfied. That hunger is useful: iodine damages key microbial molecules and kills germs. Its remaining temper also makes pure crystals and concentrated tincture irritating, requiring dilution and disposal of tested rice.

One family supplies hazards before electron gain, household materials afterward, and tools when controlled midway. Its story is written in this electron account.

\section{Beyond Giving and Taking}

Chapter 12 showed Group 1 and 2 donors; here were Group 17 takers and Group 18 bystanders. Give, take, or ignore: have we finished the table's right-hand story?

Far from it. Near the center-right live oxygen, sulfur, nitrogen, and phosphorus, neither fully surrendering electrons nor solving everything by taking. Oxygen's electron-pulling strength trails only fluorine, yet mostly works by sharing. Nitrogen's three atoms share electrons so firmly that lightning and bacteria are needed to separate them, feeding humanity through fertilizer industry. Phosphorus joins life's energy ledger. These are \textbf{life's framework elements}, forming most of your dry weight together with carbon and hydrogen. Chapter 14 visits this quartet supporting life and environment.

The halogens' remaining puzzles---shared electrons in chlorine and ionic stacking in salt---belong to chemical bonds. See Chapters 17 and 18.

\begin{tiaozhan}
Challenge, optional for the main story: if fluorine's affinity is lower than chlorine's, why is it the strongest halogen? Reactions use the whole ledger, not one entry. A typical account has at least three. (1) Splitting the halogen molecule costs energy: \chem{F_2}'s bond is unexpectedly weak, far cheaper to break than chlorine's, perhaps because crowded small atoms make lone pairs repel. (2) Electron acceptance releases affinity energy, where fluorine indeed loses slightly. (3) Water or a crystal capturing the negative ion releases energy: small, charge-concentrated \chem{F^-} binds especially strongly and releases much more. Sum all three and fluorine leads decisively. Strongest fluorine and nonhighest affinity are compatible statements about different levels of accounting. The right way to understand periodic trends is first to ask which ledger is being compared.
\end{tiaozhan}

\section*{Try It Yourself}

\begin{sikaoti}
\item Draw chlorine's 17-electron and chloride's 18-electron floors, marking incomplete and full outer shells. Explain why one appears in chemical weapons and the other in soup. Note that floor diagrams aid reasoning, not photograph real electrons.
\item Predict the color change when chlorine water enters colorless sodium iodide. Describe who gives electrons to whom and why iodine water added to sodium chloride shows nothing.
\item Bartlett compared xenon's ionization energy with oxygen's. If a substance of unknown inertness has ionization energy nearly equal to krypton's, what could you predict? Outline a test without performing it.
\item A 5-gram spoonful contains about $5\times 10^{22}$ chloride ions, each about $3.6\times 10^{-10}$ meters across. How long would they stretch in one line? Compare with the Earth--Moon distance, about $3.8\times 10^{8}$ meters.
\item Restate the reason for checking an argon balloon's identity in energy-account language: why are helium party balloons safer than hydrogen ones? Consider electron removal and addition, then hydrogen's account.
\item Draw the tungsten--halogen material cycle, with boxes for filament tungsten, evaporated tungsten, and tungsten iodide, and arrows labeled high or low temperature. Note that boxes and arrows are human representations.
\end{sikaoti}
